\chapter{Routers, cascades, and verifiers}
\label{ch:routers}

The decision layer sits in a long line of ideas for spending less compute without losing
quality. Knowing where they succeed, and where they usually fail, keeps our claims honest.

\section{Routers}

A \term{router} sends each query to the cheapest model likely to handle it. RouteLLM trains
strong-vs-weak routers from preference data and reports large savings at matched
quality~\citep{routellm}. The cautionary result: benchmarks found many routers barely beat a
trivial ``mix models at random in the right proportion'' baseline, and routers trained on
one distribution were near-random on another. Baselines matter.

\section{Cascades}

A \term{cascade} tries a cheap model first and escalates to an expensive one when a
reliability score is low. FrugalGPT is the canonical example~\citep{frugalgpt}. But theory
shows confidence-based deferral \emph{provably fails} when the downstream model is a
specialist, under label noise, and under distribution shift~\citep{cascade_deferral}.

\begin{figure}[h]\centering
\fitfig{%
\begin{tikzpicture}[node distance=6mm and 14mm, every node/.style={font=\scriptsize}]
  % router (top)
  \node[frontend] (q1) {query};
  \node[backend, right=of q1] (r) {router};
  \node[model, above right=4mm and 14mm of r] (weak) {cheap model};
  \node[model, below right=4mm and 14mm of r] (strong) {strong model};
  \draw[flow] (q1)--(r); \draw[flow] (r)--(weak); \draw[flow] (r)--(strong);
  \node[font=\scriptsize\bfseries, above=2mm of weak] {Router: pick before running};
\end{tikzpicture}}
\par\vspace{6mm}
\fitfig{%
\begin{tikzpicture}[node distance=6mm and 14mm, every node/.style={font=\scriptsize}]
  % cascade (separate picture, well below)
  \node[frontend] (q2) {query};
  \node[model, right=of q2] (cheap) {cheap model};
  \node[backend, right=of cheap] (gate) {confident?};
  \node[model, right=of gate] (exp) {strong model};
  \node[frontend, above=7mm of gate] (out) {answer};
  \draw[flow] (q2)--(cheap); \draw[flow] (cheap)--(gate);
  \draw[flow] (gate)-- node[flowlabel]{yes} (out);
  \draw[flow] (gate)-- node[flowlabel]{no} (exp);
  \node[font=\scriptsize\bfseries, below=2mm of gate] {Cascade: pay cheap first, then decide};
\end{tikzpicture}}
\caption{Router vs cascade. A router decides \emph{before} paying (top); a cascade pays the
cheap model first, then defers on low confidence (bottom).\datalabel{Illustrative}}
\label{fig:routercascade}
\end{figure}

\begin{keyidea}
Because Jev is so cheap, it acts more like a \emph{pre-generation router} (decide before
paying for the LLM) than a cascade rung. What decides whether it helps is not its latency
but its \emph{error rate on the specific decision}, and whether its confidence is calibrated
enough to gate on.
\end{keyidea}

\section{Verifiers and the ``information the proposer lacks'' rule}

Separating a \term{proposer} (the LLM) from a \term{scorer/verifier} helps in robotics and
web agents, but the literature is consistent about \emph{why}: it helps when the scorer has
information the proposer does not, such as execution results or environment state, not merely
because it is a second model. Imperfect verifiers also cap the gains and can be gamed at
large sample counts.

\begin{whyhere}
This sets our expectation for the honest outcome: ``cheaper decisions at equal or slightly
lower accuracy,'' not ``a strictly better agent.'' The evaluation is designed to detect
exactly that, and the arms (Chapter~\ref{ch:eval}) separate ``typed decision helped'' from
``Jev specifically helped'' from ``more compute helped.''
\end{whyhere}

\begin{checkbox}
Why is a cheap, calibrated classifier better modeled as a pre-generation router than as a
cascade rung? (Appendix~\ref{app:answers}.)
\end{checkbox}
